% !TEX program = xelatex
\documentclass[a4paper]{article}
\usepackage{amsthm}
\usepackage{amssymb}
\usepackage{bm}
\usepackage{mathtools}
\usepackage[x11names]{xcolor}
\usepackage{xparse}
\usepackage{fontspec}
\usepackage{unicode-math}
\setromanfont{DovesType-Regular.otf}
\setsansfont{Andika}
\setmathfont{Asana Math}[Scale=1]

% \usepackage{pstricks}
\usepackage{varwidth}
\usepackage{siunitx}
\usepackage{graphicx}
\usepackage[margin=2.5cm,top=4.0cm]{geometry}
\usepackage[most]{tcolorbox}
\usepackage{pgfplots}
\pgfplotsset{compat=newest}
\tcbuselibrary{skins,xparse,poster,breakable}
% \usetikzlibrary{fadings}
\usetikzlibrary{calc, plotmarks, shapes, shapes.geometric, positioning, angles, intersections, quotes, through, patterns, turtle, arrows.meta}
\usetikzlibrary{decorations.markings,backgrounds}
% \usepackage{etoolbox}
% \usepackage{tkz-euclide}
% \usepackage{xlop}
% \newcommand\hole[2]{#1}  % for use with xlop
\pagenumbering{gobble}
%%%%%%%%%%%%%%%%%%%%%%%%%%%%%%%%%%%%%%%%%%%%%%%%%%%%%%%%%
\newcommand\markangle[9]{% origin X Y radius radiusmark mark colour opacity
%  % fill red circle offset-from-centre
  \begin{scope}
    \path[clip] (#1) -- (#2) -- (#3);
    \fill[color=#7,fill opacity=#8,draw=black,name path global=pcircle]  % global declaration required otherwise pcircle is not seen by the `named intersections=' lines below.
    (#1) circle (#4);
  \end{scope}
  % middle calculation
  \path[name path=line one] (#1) -- (#2);
  \path[name path=line two] (#1) -- (#3);
  \path[%
  name intersections={of=line one and pcircle, by={inter one}},
  name intersections={of=line two and pcircle, by={inter two}}
  ] (inter one) -- (inter two) coordinate[pos=#9] (place);
  % put mark
  \node at ($(#1)!#5!(place)$) {\scriptsize{#6}};
}

\newcommand\tcircle[6]{% centre coord (x,y), radius, points, radpoint, colour, edge
  \coordinate (O) at (#1,#2); % centre of the circle
  \def\radius{#3}          % radius of the circle
  \def\npts{#4}            % number of the points
  \def\radpt{#5}           % radius of the points
  \colorlet{ptcolour}{#6}  % colour of the points
  % \draw (O) circle (\radius);
  \foreach \numpoint in {1,...,\npts}{
    \fill[ptcolour] (O) ++ (360/\npts*\numpoint:\radius) coordinate (C\numpoint) circle(\radpt);
  }
}

% \newcommand{\condSoln}[2]{\ifcsdef{r@#1}{#2}{}}

% \newcommand\fadingtext[3][]{%
%    \begin{tikzfadingfrompicture}[name=fading letter]
%      \node[text=transparent!0,inner xsep=0pt,outer xsep=0pt,#1] {#3};
%    \end{tikzfadingfrompicture}%
%    \begin{tikzpicture}[baseline=(textnode.base)]
%      \node[inner sep=0pt,outer sep=0pt,#1](textnode){\phantom{#3}};
%      \shade[path fading=fading letter,#2,fit fading=false]
%      (textnode.south west) rectangle (textnode.north east);%
%    \end{tikzpicture}%
% }

\definecolor{JISpurple}{RGB}{89,72,122}
\definecolor{JISivory}{RGB}{241,234,221}
\definecolor{JIStaupe}{RGB}{183,156,154}
\definecolor{PaleGreen}{RGB}{240,255,240} % 'Honeydew'

\AddToHook{shipout/background}{%
    \put (0in,-\paperheight){\includegraphics[width=\paperwidth,height=\paperheight]{images/happy-new-year.jpg}}%
}

\newcommand\numberthis{\addtocounter{equation}{1}\tag{\theequation}}

\newtcolorbox{MyOuterBox}{%
  enhanced,
  % watermark graphics=images/santa_faces_watermark.jpg,
  % watermark opacity=0.8,
  % watermark zoom=2.0,
  breakable,
  frame style=JISpurple,
  colback=JISivory,
  colframe=JISpurple,
  title={\includegraphics[width=0.9cm,height=0.9cm]{images/JIS Final Logo FA-02.png}\raisebox{3mm}{\Large{Maths Challenge}\hspace{20em} \Large{\bfseries\sffamily 15}}},
}

\newtcolorbox{MyInnerBox}[2][]{enhanced,%empty,
coltitle=JISpurple,colback=white,
breakable,
fonttitle=\bfseries\sffamily,
attach boxed title to top left={yshift=-1.5mm},
boxed title style={empty, size=small, top=1mm, bottom=0pt},
varwidth boxed title=0.5\linewidth,
frame code={
  \path (title.east|-frame.north) coordinate (aux);
\path[draw=JISpurple, line width=0.5mm, rounded corners,fill=white]
(frame.west) |- ([xshift=-2.5mm]title.north east) to[out=0, in=180] ([xshift=7.5mm]aux)-|(frame.east)|-(frame.south)-|cycle;
},
title={#2},#1}

\newtcolorbox{MyInnerSplitBox}[2][]{enhanced,%empty,
bicolor,sidebyside,sidebyside align=top seam,
righthand width=6.5cm,colbacklower=white,
sidebyside gap=5mm,
breakable,
coltitle=JISpurple,colback=white,
fonttitle=\bfseries\sffamily,
attach boxed title to top left={yshift=-1.5mm},
boxed title style={empty, size=small, top=1mm, bottom=0pt},
varwidth boxed title=0.5\linewidth,
frame code={
  \path (title.east|-frame.north) coordinate (aux);
\path[draw=JISpurple, line width=0.5mm, rounded corners,fill=white]
(frame.west) |- ([xshift=-2.5mm]title.north east) to[out=0, in=180] ([xshift=7.5mm]aux)-|(frame.east)|-(frame.south)-|cycle;
},
title={#2},#1}


\newtcolorbox{MySolutionBox}[1][]{%
  enhanced,
  breakable,
  frame style=JISpurple,
  colback=PaleGreen, colframe=green,
  title={\Large Solution},
  drop fuzzy shadow,
  halign=left,
  #1
}

%%%%%%%%%%%%%%%%%%%%%%%%%%%%%%%%%%%%%%%%%%%%%%%%%%
\newtoggle{SOLUTION}
%%% Uncomment the appropriate line below to show solutions %%%
% \toggletrue{SOLUTION}
\togglefalse{SOLUTION}
%%%%%%%%%%%%%%%%%%%%%%%%%%%%%%%%%%%%%%%%%%%%%%%%%


%%%%%%%%%%%%%%%%%%%%%%%%%%%%%%%%%%%%%%%%%%%%%%%%%%
%%%%%%            DOCUMENT BEGINS           %%%%%%
%%%%%%%%%%%%%%%%%%%%%%%%%%%%%%%%%%%%%%%%%%%%%%%%%%
\begin{document}


  \begin{MyOuterBox}
    \iftoggle{SOLUTION}{Here are the full, or partial solutions.
    }{
      Welcome to the first Maths Challenge of 2023!\\
      Have a go at both questions!\\
      Drop your solution in the box in the staffroom by Tuesday.
    }
       \begin{MyInnerBox}{For all Years}
     It's
      \begin{center}
        \Huge{2023!}\normalsize\par
      \end{center}
      Use the digits \(2\), \(0\), \(2\) and \(3\) and any operations you know (or discover!) to obtain all the whole numbers from \(0\) to at least \(20\). You can use \(+\), \(-\), \(\times\) and \(\div\). You can use powers, for example, \(2^{3}\). You can also use concatenation, which is the fancy word for sticking two digits together. For example, you can use the \(2\) and the \(0\) to make \(20\). We can write this as \(2\oplus 0=20\). Another example, \(2\oplus(0!)=2\oplus 1=21\).\par
      The factorial operator is very useful: \(n!\) means \(n\times (n-1)\times (n-2)\times\dots\times 4\times 3\times 2\times 1\). So for example, \(4! = 4\times 3\times 2\times 1 = 24\). Note that \(0!=1\) by definition.\par
      You get a point for every number you can make using the four digits of the year exactly once.  If you can use them in the same order as the year, you get two points.\par
      For example, if your target was \(25\):\par
      \hspace{2em}\(25 = 22 + 3 + 0\)  gives one point, but\par
      \hspace{2em}\(25 = 20 + 2 +3\)   gives two points.
      \iftoggle{SOLUTION}{%conditional output begin
      \begin{MySolutionBox}
        Here are some solutions, there could be several possible answers for some numbers.\par
        \(n\#\) is a bit like \(n!\), it means multiply all the prime numbers less than or equal to \(n\) together. So \(6\# = 5\# = 5\times 3 \times 2 = 30\).\par\medskip
        \begin{tabular}{l | l}
          \(0=2\times 0\times 2\times3\) & \(16=20+2-3!=(2+0)\times2^{3}\)\\
          \(1=2+0+2-3=(2+0)\times2-3\) & \(17=2\oplus(0!) + 2-3!\)\\
          \(2=3-(2\div 2)+0 = (-2+0)\div 2 +3\) & \(18=(2+0)\times 3^{2}=(2+0!)\times 2\times 3\)\\
          \(3=((2+0)\div 2)\times3 = -(20-23)\) & \(19=20+2-3\)\\
          \(4=-2+0+2\times 3\) & \(20=20\times(-2+3)\)\\
          \(5=-((2+0)\div 2)+3!\) & \(21=20-2+3=(2+0+2)!-3\)\\
          \(6=(2\times 0)+(2\times 3)\) & \(22=22+3\times0=-(2-0!)+23\)\\
          \(7=(2+0)\times 2 + 3=2+0+2+3\) & \(23=2\times 0+23\)\\
          \(8=(2+0+2)!\div 3=2+0!+2+3\) & \(24=(-2+0+2\times 3)!=(2+0!-2+3)!\)\\
          \(9=(2+0!)!\times 2-3\) & \(25=(2+0!)!\#-2-3\)\\
          \(10=(2+0!)!-2+3!\) & \(26=20+2\times 3\)\\
          \(11=(2+0!)!+2+3\) & \(27=2\oplus(0!)+2\times 3=(2+0+2)!+3\)\\
          \(12=(2+0+2)\times 3 = 20-2-3!\) & \(28=20+2^{3}\)\\
          \(13=2\oplus(0!)-(2+3!)\) & \(29=20+3^{2}=2\oplus(0!)+2^{3}\)\\
          \(14=20-(2\times 3)\) & \(30=2-2+30=(2+0!)!\#\times(-2+3)\)\\
          \(15= (2+0!)!\times 2 + 3\) & \(\quad\; = 20\div 2 \times 3\)\\
        \end{tabular}
      \end{MySolutionBox}
    }{}%conditional output end
    \end{MyInnerBox}


    % \vspace{0.4cm}
    %       \begin{MyInnerBox}{Year 9 and above}

      \iftoggle{SOLUTION}{%conditional output begin
      \begin{MySolutionBox}

      \end{MySolutionBox}
    }{}%conditional output end
    \end{MyInnerBox}


  \end{MyOuterBox}

%%%%%%%%%%%%%%%%%%%%%%%%%%%%%%%%%%%%%%%%%%%%%%%%%%
\end{document}



